\documentclass[11pt,a4paper]{article}

\usepackage[margin=2.5cm]{geometry}
\usepackage[T1]{fontenc}
\usepackage[utf8]{inputenc}
\usepackage{xcolor}
\usepackage{listings}
\usepackage{enumitem}
\usepackage[colorlinks=true,urlcolor=blue,linkcolor=black]{hyperref}
\usepackage{parskip}

% ---- Shell listing style ----
\definecolor{shellbg}{RGB}{245,247,250}
\definecolor{shellrule}{RGB}{220,225,232}
\lstdefinestyle{shell}{
  backgroundcolor=\color{shellbg},
  basicstyle=\ttfamily\small,
  frame=single,
  rulecolor=\color{shellrule},
  framesep=6pt,
  breaklines=true,
  columns=fullflexible,
  keepspaces=true,
  showstringspaces=false,
  xleftmargin=4pt,
  xrightmargin=4pt,
}
\lstset{style=shell}

\title{\textbf{CityVision --- Running the App with Docker}}
\author{Urban Scene Semantic Segmentation}
\date{}

\begin{document}
\maketitle

\noindent
The application runs as \textbf{two containers}: a \textbf{backend} (FastAPI
prediction API) and a \textbf{frontend} (Streamlit UI that consumes the API).
Both are started together with Docker Compose.

\section*{Prerequisites}
\begin{itemize}[leftmargin=1.5em]
  \item \textbf{Docker Desktop} installed and \emph{running} (the whale icon in
        the system tray is steady; the app shows ``Engine running'').
  \item A terminal opened at the project root:\\
        \texttt{C:\textbackslash Users\textbackslash DFILATOVA\textbackslash Desktop\textbackslash CityVision}
\end{itemize}

\noindent Check that Docker is running:
\begin{lstlisting}
docker version
\end{lstlisting}
If a \texttt{Server:} section is printed, the engine is up. If you only see the
\texttt{Client:} block plus an error, start Docker Desktop first.

\section*{Step 1 --- Start the containers}
From the project root:
\begin{lstlisting}
docker compose up -d
\end{lstlisting}
\begin{itemize}[leftmargin=1.5em]
  \item \texttt{-d} runs them in the background (detached).
  \item The \emph{first} run builds the images and can take a few minutes
        (it downloads PyTorch). Later runs are fast.
  \item If you changed the code, rebuild first:
        \texttt{docker compose up -{}-build -d}
\end{itemize}

\section*{Step 2 --- Check they are running}
\begin{lstlisting}
docker compose ps
\end{lstlisting}
You should see \textbf{two} containers, both \texttt{running}:
\begin{lstlisting}
cityvision-backend-1    running   0.0.0.0:8000->8000/tcp
cityvision-frontend-1   running   0.0.0.0:8501->8501/tcp
\end{lstlisting}

\section*{Step 3 --- Open the app}
\begin{itemize}[leftmargin=1.5em]
  \item \textbf{Application (frontend):} \url{http://localhost:8501}
  \item \textbf{API documentation (backend):} \url{http://localhost:8000/docs}
\end{itemize}
On the very first load, wait $\sim$10 seconds for the backend to load the model,
then refresh if you see a connection error.

\section*{Step 4 --- Use the app}
\begin{enumerate}[leftmargin=1.5em]
  \item In the sidebar, select an \textbf{image} from the dropdown
        (e.g.\ \texttt{frankfurt\_000000\_000294}).
  \item Click \textbf{Predict mask}.
  \item The three panels appear: \textbf{real image}, \textbf{real mask}, and
        \textbf{predicted mask}, plus the detected-class breakdown.
\end{enumerate}

\section*{Step 5 --- Stop the containers}
\begin{lstlisting}
docker compose down
\end{lstlisting}
This stops and removes both containers. The built images stay cached, so the
next start is quick.

\section*{Quick reference}
\begin{center}
\begin{tabular}{ll}
\textbf{Goal} & \textbf{Command} \\[2pt]
Start (images already built) & \texttt{docker compose up -d} \\
Start after code changes      & \texttt{docker compose up -{}-build -d} \\
Check status                  & \texttt{docker compose ps} \\
View logs                     & \texttt{docker compose logs -f} \\
Stop                          & \texttt{docker compose down} \\
\end{tabular}
\end{center}

\section*{Troubleshooting}
\begin{description}[leftmargin=0em,style=nextline]
  \item[\texttt{error during connect \dots dockerDesktopLinuxEngine}]
        Docker Desktop's engine is not running. Open Docker Desktop and wait for
        ``Engine running'', then retry.
  \item[\texttt{az} / other command ``not recognized'']
        Not related to running the app; only relevant for cloud deployment.
  \item[Frontend shows ``API not reachable'']
        The backend is still starting (model load). Wait a few seconds and
        refresh. Confirm with \texttt{docker compose ps} that
        \texttt{cityvision-backend-1} is \texttt{running}.
  \item[Port already in use]
        Another process holds port 8501 or 8000. Stop it, or change the port
        mapping in \texttt{docker-compose.yml}.
\end{description}

\end{document}
